% !TeX root = ../../Book1.tex
% 纲要标题：### 6.2 半直积
% 纲要位置：计划纲要.md，第 284 行。

\section{半直积}
\label{b1:sec:0602}

直积中来自不同因子的元素总是交换，但二面体群的旋转与反射并不满足这一条件：反射会把旋转共轭为它的逆。半直积保留了两因子给出的坐标，同时用一个自同构作用记录它们怎样交换位置。这里既要从群和作用构造新群，也要判断一个已有的群何时能作这种分解；二者的乘法公式应当相互对应。

% 本节正文按下列原纲要要点撰写。

% 纲要内容（仅作写作计划，尚非教材正文）：
%
% * 自同构作用与外半直积
% * 内半直积判别
% * 二面体群、仿射群与具体计算
%

\subsection{自同构作用与外半直积}
\label{b1:subsec:0602:01}

设想在一个群 \(G\) 中，\(N\) 是正规子群，\(H\) 是子群，每个元素都能唯一写成 \(nh\)，其中 \(n\in N\)、\(h\in H\)。相乘时，要把第二个元素的 \(N\) 部分移过第一个元素的 \(H\) 部分。插入 \(h^{-1}h=1\)，得到
\[
(nh)(n'h')=n(hn'h^{-1})hh'.
\]
正规性使 \(hn'h^{-1}\) 仍在 \(N\) 中；把这个共轭作用记为 \(\alpha_h(n')\)，乘积的坐标就必须是
\[
(n,h)(n',h')=\bigl(n\cdot\alpha_h(n'),hh'\bigr).
\]
下面反过来，只给定两个群和一个自同构作用，用这一公式构造群。

\begin{definition}\label{b1:def:semidirect-product}
设 \(N,H\) 为群，给定群同态 \(\alpha:H\rightarrow\operatorname{Aut}(N)\)，并记 \(\alpha_h=\alpha(h)\)。也就是说，每个 \(\alpha_h:N\rightarrow N\) 是自同构，且 \(\alpha_1=\operatorname{id}_N\)、\(\alpha_{hh'}=\alpha_h\circ\alpha_{h'}\)。在集合 \(N\times H\) 上规定乘法
\[
(n,h)(n',h')=\bigl(n\cdot\alpha_h(n'),hh'\bigr)
\]
称由此得到的群为 \(N,H\) 相对于 \(\alpha\) 的\term[wai-banzhiji]{外半直积}[external semidirect product]，记为 \(N\rtimes_\alpha H\)；作用明确时简记 \(N\rtimes H\)。
\symidx{\(N\rtimes_\alpha H\)：外半直积}
\end{definition}

这里先验证它确为群。三个元素相乘时，左结合的第一坐标为
\[
n\cdot\alpha_h(n')\alpha_{hh'}(n''),
\]
右结合的第一坐标为
\[
n\cdot\alpha_h\bigl(n'\cdot\alpha_{h'}(n'')\bigr)
=n\cdot\alpha_h(n')(\alpha_h\alpha_{h'})(n''),
\]
两者相等，使用了 \(\alpha_h\) 保持乘法和 \(\alpha_{hh'}=\alpha_h\alpha_{h'}\)；第二坐标的结合律来自 \(H\)。由 \(\alpha_1=\operatorname{id}_N\) 及 \(\alpha_h(1)=1\)，\((1,1)\) 是单位。求 \((n,h)\) 的右逆时，第二坐标必须是 \(h^{-1}\)，而第一坐标 \(x\) 须满足
\[
(n,h)(x,h^{-1})=(1,1)
\quad\Longleftrightarrow\quad
n\cdot\alpha_h(x)=1.
\]
因此 \(\alpha_h(x)=n^{-1}\)，即 \(x=\alpha_h^{-1}(n^{-1})\)。因为 \(\alpha\) 是群同态，
\[
\alpha_h\alpha_{h^{-1}}=\alpha_{h^{-1}}\alpha_h
=\alpha_1=\operatorname{id}_N,
\]
所以 \(\alpha_h^{-1}=\alpha_{h^{-1}}\)，得到
\[
(n,h)^{-1}=\bigl(\alpha_{h^{-1}}(n^{-1}),h^{-1}\bigr),
\]
这里把逆自同构写成 \(\alpha_{h^{-1}}\)，用到的是 \(\alpha\) 保持乘法，而不只是每个 \(\alpha_h\) 都为自同构。上述元素也为左逆，因为
\[
\bigl(\alpha_{h^{-1}}(n^{-1}),h^{-1}\bigr)(n,h)
=\bigl(\alpha_{h^{-1}}(n^{-1})\alpha_{h^{-1}}(n),1\bigr)
=(1,1).
\]

嵌入 \(n\mapsto(n,1)\)、\(h\mapsto(1,h)\) 把 \(N,H\) 识别为子群；投影 \((n,h)\mapsto h\) 是满同态，核为 \(N\times\{1\}\)，故这个副本的 \(N\) 正规。两种嵌入满足
\begin{equation}\label{b1:eq:semidirect-conjugation}
(1,h)(n,1)(1,h)^{-1}=\bigl(\alpha_h(n),1\bigr).
\end{equation}
此外，每个元素唯一写为 \((n,1)(1,h)\)，所以一个同态在这两个坐标子群上的限制已经决定它在全群上的值。给定到同一个群 \(K\) 的同态 \(u:N\rightarrow K\) 和 \(v:H\rightarrow K\)，除了各自保持因子内部的乘法，还必须保留\eqref{b1:eq:semidirect-conjugation}中的共轭关系。这正是下面映射性质里的相容条件；证明将说明，它也是存在所需同态的充分条件。

\begin{proposition}[半直积的映射性质]\label{b1:prop:semidirect-universal}
设 \(N,H,K\) 为群，\(\alpha:H\rightarrow\operatorname{Aut}(N)\) 为群同态，并由此构造半直积 \(N\rtimes_\alpha H\)。给定群同态 \(u:N\rightarrow K\)、\(v:H\rightarrow K\)，存在群同态
\(F:N\rtimes_\alpha H\rightarrow K\) 限制为 \(u,v\)，当且仅当
\[
v(h)u(n)v(h)^{-1}=u\bigl(\alpha_h(n)\bigr)
\quad(n\in N,\ h\in H).
\]
存在时唯一，并且 \(F(n,h)=u(n)v(h)\)。
\end{proposition}
\begin{proof}
若 \(F\) 是所求同态，对\eqref{b1:eq:semidirect-conjugation}两边应用 \(F\)，并代入 \(F(n,1)=u(n)\)、\(F(1,h)=v(h)\)，就得到
\(v(h)u(n)v(h)^{-1}=u\bigl(\alpha_h(n)\bigr)\)，这证明必要性。

反过来，满足此条件时定义 \(F(n,h)=u(n)v(h)\)。计算
\[
\begin{aligned}
u(n)v(h)u(n')v(h')
&=u(n)u\bigl(\alpha_h(n')\bigr)v(h)v(h')\\
&=u\bigl(n\cdot\alpha_h(n')\bigr)v(hh'),
\end{aligned}
\]
故 \(F(n,h)F(n',h')=F\bigl(n\cdot\alpha_h(n'),hh'\bigr)\)，正是保持半直积乘法的条件。又因 \(u(1)=v(1)=1\)，此映射在两个坐标子群上的限制确为 \(u,v\)。

最后，任何满足限制条件的同态都必须把 \((n,h)=(n,1)(1,h)\) 送到 \(u(n)v(h)\)，所以唯一。
\end{proof}

\subsection{内半直积判别}
\label{b1:subsec:0602:02}

\begin{theorem}[内半直积判别]\label{b1:thm:internal-semidirect}
设 \(G\) 为群，\(N\trianglelefteq G\)、\(H\leq G\)，并且 \(N\cap H=\{1_G\}\)、\(NH=G\)。共轭给出同态
\[
\alpha:H\longrightarrow\operatorname{Aut}(N),\qquad
\alpha_h(n)=hnh^{-1}.
\]
映射 \(N\rtimes_\alpha H\rightarrow G\)、\((n,h)\mapsto nh\) 是同构。
\end{theorem}
\begin{proof}
正规性保证 \(\alpha_h\) 取值于 \(N\)，其逆为 \(\alpha_{h^{-1}}\)，且共轭保持乘法；\(\alpha_{hh'}=\alpha_h\alpha_{h'}\) 由直接计算成立。

条件 \(NH=G\) 保证每个 \(g\in G\) 至少有一种表达 \(g=nh\)，所以 \((n,h)\mapsto nh\) 满射。条件 \(N\cap H=\{1_G\}\) 则保证坐标唯一：若 \(nh=n'h'\)，则
\[
n'^{-1}n=h'h^{-1}\in N\cap H=\{1_G\},
\]
故 \(n=n'\)、\(h=h'\)，映射单射。

最后计算乘法：
\[
(nh)(n'h')=n(hn'h^{-1})hh'=n\cdot\alpha_h(n')hh'.
\]
因此这个双射也保持半直积乘法，是群同构。
\end{proof}

设 \(G\) 为群，\(N\trianglelefteq G\)、\(H\leq G\)。如果 \(N\cap H=\{1_G\}\) 且 \(NH=G\)，则将 \(G\) 称为 \(N,H\) 的\term[nei-banzhiji]{内半直积}[internal semidirect product]，并将 \(H\) 称为 \(N\) 的一个\term[buqun]{补群}[complement]。这一分解不仅说 \(N,H\) 生成 \(G\)，还为每个元素给出唯一的 \(N\)-坐标与 \(H\)-坐标。

\ProseParagraphBreak
补群是一种附加选择。仅有 \(N\trianglelefteq G\) 并不保证补群存在，即使存在也未必唯一；下面仿射群的例子将展示不同的补群。补群的存在与下一节定义的扩张分裂相联系，\cref{b1:thm:split-extension-semidirect}将给出精确的等价条件。

\ProseParagraphBreak
直积是作用平凡的半直积：当 \(\alpha_h(n)=n\) 时，乘法恢复为逐坐标乘法。反过来，若此内分解中的 \(H\) 也正规，则满足\cref{b1:prop:internal-direct-product}的全部条件；该命题证明中的计算给出 \(hnh^{-1}n^{-1}\in N\cap H=\{1\}\)，故 \(\alpha_h(n)=n\)，作用便平凡。

\ProseParagraphBreak
外半直积从群和作用制造新群，内半直积则在一个既有群中验证能否如此分解。符号 \(N\rtimes H\) 仅省略了已经说明的作用，不表示同一对抽象群只有一种半直积；例如 \(C_3\) 与 \(C_2\) 的平凡作用给出 \(C_6\)，取逆作用给出 \(S_3\)。

\subsection{二面体群、仿射群与具体计算}
\label{b1:subsec:0602:03}

在 \(D_{2n}\) 中，取旋转子群 \(N=\langle r\rangle\cong C_n\) 和反射子群 \(H=\langle s\rangle\cong C_2\)，其中 \(n\geq3\)。\cref{b1:prop:dihedral-normal-form}给出每个对称的唯一表达 \(r^is^j\) 及关系 \(srs^{-1}=r^{-1}\)。因 \(\langle r\rangle\) 被生成元 \(r,s\) 的共轭保持，\(N\trianglelefteq D_{2n}\)；唯一表达又给出 \(N\cap H=\{1\}\) 和 \(NH=D_{2n}\)。因此可应用\cref{b1:thm:internal-semidirect}，其中 \(s\) 的共轭作用为取逆，得到
\[
D_{2n}\cong C_n\rtimes C_2,\qquad
(i,j)(k,\ell)=\bigl(i+(-1)^jk,j+\ell\bigr),
\]
第一坐标模 \(n\)，第二坐标模 \(2\)。例如 \((i,1)^2=(0,0)\)，所以每个反射阶为 \(2\)；而 \((1,1)(1,0)=(0,1)\)、\((1,0)(1,1)=(2,1)\)，在 \(n\geq3\) 时不同，明确显示非交换性。

\ProseParagraphBreak
另一个例子来自实直线上的仿射变换 \(x\mapsto ax+b\)，其中 \(a\in\mathbb R^\times\)、\(b\in\mathbb R\)。先执行 \(x\mapsto a'x+b'\)，再执行 \(x\mapsto ax+b\)，得到
\[
x\longmapsto aa'x+(b+ab').
\]
所以用 \((b,a)\) 记录变换时，乘法为
\[
(b,a)(b',a')=(b+ab',aa'),
\]
从而\term[fangshe-qun]{仿射群}[affine group]为
\((\mathbb R,+)\rtimes(\mathbb R^\times,\cdot)\)，作用为数乘。它也可写成矩阵群
\[
\Biggl\{\begin{pmatrix}a&b\\0&1\end{pmatrix}
\ \Bigg|\ a\neq0,\ b\in\mathbb R\Biggr\}.
\]
平移组成正规子群 \(T=\{\,(b,1)\mid b\in\mathbb R\,\}\)，固定 \(0\) 的伸缩组成子群 \(L_0=\{\,(0,a)\mid a\in\mathbb R^\times\,\}\)。由 \((b,a)=(b,1)(0,a)\) 及 \(T\cap L_0=\{(0,1)\}\)，\(L_0\) 是 \(T\) 的一个补群。

\ProseParagraphBreak
若把原点改为 \(c\in\mathbb R\)，相应的伸缩为
\[
f_{a,c}(x)=a(x-c)+c=ax+(1-a)c,
\qquad a\in\mathbb R^\times.
\]
记这些变换组成的集合为 \(L_c\)，并令 \(\tau_c(x)=x+c\)、\(f_a(x)=ax\)。直接计算得
\[
(\tau_c\circ f_a\circ\tau_c^{-1})(x)=a(x-c)+c,
\qquad L_c=\tau_c L_0\tau_c^{-1}.
\]
由于 \(\tau_c\in T\) 且 \(T\) 正规，共轭把 \(T\) 保持为自身，因而也把补群 \(L_0\) 变成补群 \(L_c\)。不同的 \(c\) 给出不同的补群：\(L_c\) 中斜率为 \(-1\) 的变换是 \(x\mapsto-x+2c\)，它确定了 \(c\)。所以换原点改变的是补群的选择，而整个仿射群保持不变。

\ProseParagraphBreak
\cref{b1:thm:groups-order-pq}中的乘法现在可写成
\(C_q\rtimes C_p\)，其中 \(C_p\) 的生成元对 \(C_q\) 的加法剩余类作用为乘以 \(t\)。
